\section{Methodology}

{\method} aims to generate large-scale 3D driving scenes within a unified framework addressing controllability, fidelity, and scalability.
As shown in Fig.~\ref{fig:pipeline}, it consists of three main components:
\textbf{1) Multi-Granular Controllability} (Sec.~\ref{sec_control}), which integrates high-level user intent with low-level geometric constraints for flexible scene specification;
\textbf{2) Joint Occupancy, Image, and Video Generation} (Sec.~\ref{sec_generation}), which employs conditioned diffusion models to synthesize 3D voxel occupancy, multi-view images, and temporally coherent videos with 3D-aware guidance; and
\textbf{3) Large-Scale Scene Extrapolation and Reconstruction} (Sec.~\ref{sec_extrapolate}), which extends scenes via consistency-aware outpainting and lifts them into 3DGS representations for downstream simulation and exploration.


\subsection{Multi-Granular Controllability}
\label{sec_control}
{\method} supports dual-mode scene control through: 1) high-level textual prompts, which are enriched by LLMs and converted into structured layouts via a text-to-layout generation model (illustrated in Fig.~\ref{fig:text2layout}); and 2) direct low-level geometric control for precise spatial specification.
This hybrid approach enables both intuitive creative expression and exacting scene customization.

\begin{figure}[t]
    \centering
    \includegraphics[width=\textwidth]{src/text2layout.pdf}
    \vspace{-1.5em}
    \caption{\textbf{Pipeline of textual description enrichment and scene-graph to layout generation}: (a) Input prompts are enriched using RAG-augmented LLMs to produce structured scene descriptions; (b) Spatial relationships are converted into a scene graph and encoded with a graph network, followed by conditional diffusion that denoises object boxes and lane polylines into the final layouts.}
    \vspace{-1em}
\label{fig:text2layout}
\end{figure}






\paragraph{Text Description Enrichment.}
Given a coarse user-provided textual prompt $\mathcal{T}_\mathcal{P}$, we first enrich it into a comprehensive scene description $\mathcal{D} = \{\mathcal{S}, \mathcal{O}, \mathcal{B}, \mathcal{L}\}$, comprising: scene style $\mathcal{S}$ (weather, lighting, environment), foreground objects $\mathcal{O}$ (semantics, spatial attributes, and appearance), background elements $\mathcal{B}$ (semantics and visual characteristics), and textual scene-graph layout $\mathcal{L}$, representing spatial relationships among scene entities.
The structured description $\mathcal{D}$ is generated as:
\begin{equation}
    \mathcal{D} = \mathcal{G}_\text{description}\big(\mathcal{T}_\mathcal{P}, \mathrm{RAG}(\mathcal{T}_\mathcal{P}, \mathcal{M})\big)
\end{equation}
where $\mathcal{M} = \{m_i\}_{i=1}^{N}$ denotes the scene description memory. Each entity $m_i$ is automatically constructed using one of the collected scene datasets by: 
1) extracting $\{\mathcal{S}, \mathcal{O}, \mathcal{B}\}$ using VLMs on scene images; and 
2) converting spatial annotations (object boxes and road lanes) into textual scene-graph layout $\mathcal{L}$. 
As shown in Fig.~\ref{fig:text2layout}, the Retrieval-Augmented Generation (RAG) module retrieves relevant descriptions similar to $\mathcal{T}_\mathcal{P}$ from the memory bank $\mathcal{M}$, which are then composed into a detailed, user-intended scene description by an LLM-based generator $\mathcal{G}_{\text{description}}$.

This pipeline leverages RAG for few-shot retrieval and composition when processing brief user prompts, enabling flexible and context-aware scene synthesis. The memory bank $\mathcal{M}$ is designed to be extensible, allowing seamless integration of new datasets to support a broader variety of scene styles. Additional examples of generated scene descriptions are provided in the appendix.

\paragraph{Textual Scene-Graph to Layout Generation.}
Given the textual layout $\mathcal{L}$, we follow prior works~\cite{commonscenes, echoscene} and translate it into a spatial layout map via a scene-graph–to–layout pipeline (Fig.~\ref{fig:text2layout}).
First, we construct a scene graph $\mathcal{G}=(\mathcal{V}, \mathcal{E})$, where nodes $\mathcal{V} = \{v_i\}_{i=1}^M$ represent $M$ scene entities (e.g., \textit{cars}, \textit{pedestrians}, \textit{road lanes}) and edges $\mathcal{E} = \{e_{i \rightarrow j}|i,j \in\{1,...,M\}\}$ represent spatial relations (e.g., \textit{front of}, \textit{on top of}). 
Each node and edge is then embedded by concatenating semantic features $s_i$, $s_{i \rightarrow j}$ (extracted via a text encoder $\mathcal{E}_{\text{text}}$) with learnable geometric features $g_i$, $g_{i \rightarrow j}$, resulting in node embeddings $\mathbf{v}_i = \text{Concat}(s_i, g_i)$ and edge embeddings $\mathbf{e}_{i \rightarrow j} = \text{Concat}(s_{i \rightarrow j}, g_{i \rightarrow j})$.

The graph embeddings are refined using a graph convolutional network, which propagates contextual information $\mathbf{e}_{i \rightarrow j}$ across the graph and updates each node embedding $\mathbf{v}_i$ via neighborhood aggregation.
Finally, layout generation is formulated as a conditional diffusion process: each object layout is initialized as a noisy 7-D vector $b_i \in \mathbb{R}^7$ (representing box center, dimensions, and orientation), while each road lane begins as a set of $N$ noisy 2D points $p_i \in \mathbb{R}^{N \times 2}$, with denoising process is conditioned on the corresponding node embeddings $\mathbf{v}_i$ to produce geometrically coherent placements.

\paragraph{Low-Level Conditional Encoding.}
We encode fine-grained conditions (such as user-provided or model-generated layout maps and 3D bounding boxes) into embeddings to enable precise geometric control. As illustrated in Fig.~\ref{fig:pipeline}, the 2D layout maps are processed by a ConvNet ($\mathcal{E}_\textit{layout}$) to extract layout embeddings $\mathbf{e}_\textit{layout}$, while 3D box embeddings $\mathbf{e}_\textit{box}$ are obtained via MLPs ($\mathcal{E}_{\textit{box}}$), which fuse object class and spatial coordinate features.
To further enhance geometric alignment, we project both the scene layout and 3D boxes into the camera view to generate perspective maps, which are encoded by another ConvNet ($\mathcal{E}_\textit{persp.}$) to capture spatial constraints from the image plane.
Additionally, high-level scene descriptions $\mathcal{D}$ are embedded via a T5 encoder ($\mathcal{E}_{\textit{text}}$), providing rich semantic cues for controllable generation through the resulting text embeddings $\mathbf{e}_\textit{text}$.


\subsection{Joint Occupancy, Image, and Video Generation}
\label{sec_generation}
We adopt a joint 3D-to-2D generation hierarchy that first models scene geometry via occupancy diffusion, followed by image synthesis guided by occupancy-rendered maps to ensure geometric consistency. The pipeline is further extended with a temporal diffusion module for video generation, producing smooth motion and cross-view temporal coherence.

\paragraph{Occupancy Generation via Triplane Diffusion.}
We adopt a triplane representation~\cite{Triplane} to encode 3D occupancy fields with high geometric fidelity. Given an occupancy volume $\mathbf{o} \in \mathbb{R}^{X \times Y \times Z}$, a triplane encoder compresses it into three orthogonal latent planes $\mathbf{h} = \{\mathbf{h}^{xy}, \mathbf{h}^{xz}, \mathbf{h}^{yz}\}$ with spatial downsampling. To mitigate information loss due to reduced resolution, we propose a novel triplane deformable attention mechanism that aggregates richer features for a query point $\mathbf{q} = (x, y, z)$ as:
\vspace{-1mm}
\begin{equation}
\mathbf{F}_\mathbf{q}(x, y, z) = \sum_{\mathcal{P} \in \{xy, xz, yz\}} 
\sum_{k=1}^K 
\sigma\big(\mathbf{W}_\omega^\mathcal{P} \cdot \text{PE}(x, y, z)\big)_k 
\cdot \mathbf{h}^\mathcal{P}\Big(
\text{proj}_\mathcal{P}(x, y, z) + 
\Delta p_k^\mathcal{P}
\Big)
\end{equation}
where $K$ is the number of sampling points, $\text{PE}(\cdot): \mathbb{R}^3 \rightarrow \mathbb{R}^D$ denotes positional encoding, and $\mathbf{W}_\omega^\mathcal{P} \in \mathbb{R}^{K \times D}$ generates attention weights with the softmax function $\sigma(\cdot)$. The projection function $\text{proj}_\mathcal{P}$ maps 3D coordinates to 2D planes (e.g., $\text{proj}_{xy}(x, y, z) = (x, y)$), and the learnable offset $\Delta p_k^\mathcal{P} = \mathbf{W}_o^\mathcal{P}[k] \cdot \text{PE}(x, y, z) \in \mathbb{R}^2$ uses weights $\mathbf{W}_o^\mathcal{P} \in \mathbb{R}^{2 \times D}$ to shift sampling positions for better feature alignment. Then the triplane-VAE decoder reconstructs the 3D occupancy field from the aggregated features $\textbf{F}_\textbf{q}$.

Building on the latent triplane representation $\mathbf{h}$, we introduce a conditional diffusion model $\epsilon_{\theta}^\textit{occ}$ that synthesizes novel triplanes through iterative denoising. At each timestep $t$, the model refines a noisy triplane $\mathbf{h}_t$ toward the clean target $\mathbf{h}_0$ using two complementary conditioning strategies: 1) additive spatial conditioning with the layout embedding $\mathbf{e}_\text{layout}$; and 2) cross-attention-based conditioning with $\mathcal{C} = \text{Concat}(\mathbf{e}_\text{box}, \mathbf{e}_\text{text})$, integrating geometric and semantic constraints. The model is trained to predict the added noise $\epsilon$ using the denoising objective: 
$ \mathcal{L}_\textit{diff}^\textit{occ} = \mathbb{E}_{t,\mathbf{h}_0,\epsilon} \left[ \| \epsilon - \epsilon_{\theta}^\textit{occ}(\mathbf{h}_t, t, \mathbf{e}_\text{layout}, \mathcal{C}) \|_2^2 \right] $.


\paragraph{Image Generation with 3D Geometry Guidance.}
After obtaining the 3D occupancy, we convert voxels into 3D Gaussian primitives parameterized by voxel coordinates, semantics, and opacity, which are rendered into semantic and depth maps via tile-based rasterization~\cite{3DGS}. 
To incorporate object-level geometry, we first generate normalized 3D coordinates for the entire scene and extract object-specific regions based on bounding boxes. The corresponding coordinates are encoded into object positional embeddings $\mathbf{e}_\text{pos}$, providing fine-grained geometric guidance.
The semantic, depth, and layout (or perspective) maps are processed by ConvNets and fused with $\mathbf{e}_\text{pos}$ to form the final geometric embedding $\mathbf{e}_\text{geo}$. This embedding is combined with noisy image latents to achieve pixel-aligned geometric conditioning.
The image diffusion model $\epsilon_{\theta}^{\text{img}}$ further leverages cross-attention with conditions $\mathcal{C}$ (text, camera, and box embeddings) for appearance control. The model is trained via:
$\mathcal{L}_{\text{diff}}^{\text{img}} = \mathbb{E}_{t,\mathbf{x}_0,\epsilon} \left[ \| \epsilon - \epsilon_{\theta}^{\text{img}}(\mathbf{x}_t, t, \mathbf{e}_\text{geo}, \mathcal{C}) \|_2^2 \right]$.


\paragraph{Video Generation with Motion-Aware Diffusion.}
After obtaining multi-view images, we extend the diffusion framework to synthesize temporally coherent videos conditioned on motion cues. The generated images from preceding clips serve as \textit{reference frames} to guide the denoising of subsequent noisy latents $\mathbf{x}_t$. The diffusion model $\epsilon_{\theta}^{\text{vid}}$ takes both $\mathbf{x}_t$ and encoded reference features $\mathbf{F}_{\text{ref}}$, concatenated along the temporal dimension, and applies a temporal self-attention layer to capture motion correspondences, with the relative ego poses $\mathbf{P}_{\text{rel}}$ also encoded for motion-aware conditioning. 

Only the temporal attention layers are fine-tuned from the pre-trained image diffusion model, enabling efficient transfer from spatial to temporal domains. The training objective follows the denoising formulation:
$\mathcal{L}_{\text{diff}}^{\text{vid}} = \mathbb{E}_{t,\mathbf{x}_0,\epsilon} [ \| \epsilon - \epsilon_{\theta}^{\text{vid}}(\mathbf{x}_t, t, \mathbf{F}_{\text{ref}}, \mathbf{P}_{\text{rel}}, \mathcal{C}) \|_2^2 ]$.
During inference, an \textit{autoregressive} strategy is employed for streaming video generation, where previously generated frames are reused as motion references to ensure smooth transitions and temporal coherence across clips.


\subsection{Large-Scale Scene Extrapolation and Reconstruction}
\label{sec_extrapolate}
Building on single-chunk generation, we propose a progressive extrapolation approach that coherently expands occupancy and images across multiple chunks, maintaining geometric and visual consistency with the generated multi-view videos for downstream applications.

\paragraph{Geometry-Consistent Scene Outpainting.}
We extend the occupancy field via triplane extrapolation~\cite{BlockFusion}, which decomposes the task into extrapolating three orthogonal 2D planes, as illustrated in Fig.~\ref{fig:outpaint}. The core idea is to generate a new latent plane $\mathbf{h}_0^{\text{new}}$ by synchronizing its denoising process with the forward diffusion of a known reference plane $\mathbf{h}_0^{\text{ref}}$, guided by an overlap mask $\mathbf{M}$.
Specifically, at each denoising step $t$, the new latent is updated as:
\begin{equation}
\mathbf{h}^{\text{new}}_{t-1} \gets \left( \sqrt{\bar{\alpha}_t}\mathbf{h}^{\text{ref}}_0 + \sqrt{1 - \bar{\alpha}_t} \boldsymbol{\epsilon} \right) \odot \mathbf{M} + \epsilon_{\theta}^{\textit{occ}}(\mathbf{h}^{\text{new}}_t, t) \odot (1 - \mathbf{M})
\end{equation}
\begin{wrapfigure}{r}{0.56\textwidth}
    \centering
    \includegraphics[width=1.0\linewidth]{src/outpainting.pdf}
    \vspace{-0.5cm}
    \caption{\textbf{Illustration of (a) consistency-aware outpainting:} (b) Occupancy triplane extrapolation is decomposed into three 2D plane extensions guided by overlapped regions; (c) Image extrapolation is performed via diffusion conditioned on images and camera parameters.}
    \label{fig:outpaint}
    \vspace{-1.0cm}
\end{wrapfigure}

where $\boldsymbol{\epsilon} \sim \mathcal{N}(\mathbf{0}, \mathbf{I})$ and $\bar{\alpha}_t$ is determined by the noise scheduler at timestep $t$.
This method preserves structural consistency in overlapping regions while plausibly extending reference content into unseen areas, yielding coherent and geometry-consistent scene extensions.

\paragraph{Visual-Coherent Image Extrapolation.}
Beyond occupancy outpainting, we further extrapolate the visual field for synchronized image generation.
To maintain visual coherence between the reference image $\mathbf{x}_0^{\text{ref}}$ and the new view $\mathbf{x}_0^{\text{new}}$, a naive approach warps $\mathbf{x}_0^{\text{ref}}$ using the camera pose $(R, T)$ and applies image inpainting (Fig.~\ref{fig:outpaint}).
However, using only warped images as conditions is inadequate.
To address this, we fine-tune the diffusion model $\epsilon_{\theta}^{\text{img}}$ with explicit conditioning on $\mathbf{x}_0^{\text{ref}}$ and camera embeddings $\mathbf{e}(R, T)$.
Concretely, $\mathbf{x}_0^{\text{ref}}$ is concatenated with the novel latent $\mathbf{x}_t^{\text{new}}$, and $\mathbf{e}(R, T)$ is incorporated via cross-attention, enabling view-consistent extrapolation with photorealistic visual results.